% GENERATED FILE. Edit canonical lessons, then run npm run generate:books.
% canonical-source-hash: fnv1a64:1794eae4bbbed60d
% canonical-lessons: TA-C192-kavalar, TA-C192-ottunar, TA-C192-tapalkarar, TA-C192-viyapari, TA-C192-caka-uliyar

\chapter{Police officer, Driver, Postman, Trader, Colleague}
\label{ch:ta-police-officer-driver-postman-trader-colleague}
\hlchaptermodality{192}

\begin{chapteropening}
\textbf{By the end of this chapter:} \emph{I can say a police officer, a driver, a postman, a trader and a colleague.}

Complete the last lesson of chapter 192: I can say a police officer, a driver, a postman, a trader and a colleague.
\end{chapteropening}

\section[\texorpdfstring{\ta{காவலர்} (kāvalar)}{காவலர் (kāvalar)}]{\ta{காவலர்} (kāvalar) — a police officer}

\label{lesson:TA-C192-kavalar}



\begin{quote}
Before the new one: say the Tamil for a tailor, then the Tamil for small change.
\end{quote}

\subsection*{You'll want to know: \ta{காவலர்}}
\textbf{\ta{காவலர்}} — \emph{kāvalar} — \textquotedblleft{}a police officer\textquotedblright{}.

\begin{grammarlens}[title={What you've built}]
One more word for this chapter.
\end{grammarlens}

\subsection*{Your turn}
\begin{itemize}
\raggedright
  \item \emph{Say it:} \emph{kāvalar}
  \item \emph{Say it:} \emph{kāvalar}, once more
  \item \emph{Say it:} say \emph{kāvalar}
  \item \emph{Recall:} say the Tamil for a tailor, then the Tamil for small change, then say \emph{kāvalar} again
\end{itemize}

\begin{tcolorbox}[breakable,colback=teal!4,colframe=teal!35!black,title={Before you move on}]
What does \ta{காவலர்} mean? (\textquotedblleft{}A police officer\textquotedblright{}.) Say it once more.
\end{tcolorbox}
\section[\texorpdfstring{\ta{ஓட்டுநர்} (ōṭṭunar)}{ஓட்டுநர் (ōṭṭunar)}]{\ta{ஓட்டுநர்} (ōṭṭunar) — a driver}

\label{lesson:TA-C192-ottunar}



\begin{quote}
Before the new one: say the Tamil for small change, then the Tamil for a police officer.
\end{quote}

\subsection*{You'll want to know: \ta{ஓட்டுநர்}}
\textbf{\ta{ஓட்டுநர்}} — \emph{ōṭṭunar} — \textquotedblleft{}a driver\textquotedblright{}.

\begin{grammarlens}[title={What you've built}]
One more word for this chapter.
\end{grammarlens}

\subsection*{Your turn}
\begin{itemize}
\raggedright
  \item \emph{Say it:} \emph{ōṭṭunar}
  \item \emph{Say it:} \emph{ōṭṭunar}, once more
  \item \emph{Say it:} say \emph{ōṭṭunar}
  \item \emph{Recall:} say the Tamil for small change, then the Tamil for a police officer, then say \emph{ōṭṭunar} again
\end{itemize}

\begin{tcolorbox}[breakable,colback=teal!4,colframe=teal!35!black,title={Before you move on}]
What does \ta{ஓட்டுநர்} mean? (\textquotedblleft{}A driver\textquotedblright{}.) Say it once more.
\end{tcolorbox}
\section[\texorpdfstring{\ta{தபால்காரர்} (tapālkārar)}{தபால்காரர் (tapālkārar)}]{\ta{தபால்காரர்} (tapālkārar) — a postman}

\label{lesson:TA-C192-tapalkarar}



\begin{quote}
Before the new one: say the Tamil for a police officer, then the Tamil for a driver.
\end{quote}

\subsection*{You'll want to know: \ta{தபால்காரர்}}
\textbf{\ta{தபால்காரர்}} — \emph{tapālkārar} — \textquotedblleft{}a postman\textquotedblright{}.

\begin{grammarlens}[title={What you've built}]
One more word for this chapter.
\end{grammarlens}

\subsection*{Your turn}
\begin{itemize}
\raggedright
  \item \emph{Say it:} \emph{tapālkārar}
  \item \emph{Say it:} \emph{tapālkārar}, once more
  \item \emph{Say it:} say \emph{tapālkārar}
  \item \emph{Recall:} say the Tamil for a police officer, then the Tamil for a driver, then say \emph{tapālkārar} again
\end{itemize}

\begin{tcolorbox}[breakable,colback=teal!4,colframe=teal!35!black,title={Before you move on}]
What does \ta{தபால்காரர்} mean? (\textquotedblleft{}A postman\textquotedblright{}.) Say it once more.
\end{tcolorbox}
\section[\texorpdfstring{\ta{வியாபாரி} (viyāpāri)}{வியாபாரி (viyāpāri)}]{\ta{வியாபாரி} (viyāpāri) — a trader, a merchant}

\label{lesson:TA-C192-viyapari}



\begin{quote}
Before the new one: say the Tamil for a driver, then the Tamil for a postman.
\end{quote}

\subsection*{You'll want to know: \ta{வியாபாரி}}
\textbf{\ta{வியாபாரி}} — \emph{viyāpāri} — \textquotedblleft{}a trader, a merchant\textquotedblright{}.

\begin{grammarlens}[title={What you've built}]
One more word for this chapter.
\end{grammarlens}

\subsection*{Your turn}
\begin{itemize}
\raggedright
  \item \emph{Say it:} \emph{viyāpāri}
  \item \emph{Say it:} \emph{viyāpāri}, once more
  \item \emph{Say it:} say \emph{viyāpāri}
  \item \emph{Recall:} say the Tamil for a driver, then the Tamil for a postman, then say \emph{viyāpāri} again
\end{itemize}

\begin{tcolorbox}[breakable,colback=teal!4,colframe=teal!35!black,title={Before you move on}]
What does \ta{வியாபாரி} mean? (\textquotedblleft{}A trader, a merchant\textquotedblright{}.) Say it once more.
\end{tcolorbox}
\section[\texorpdfstring{\ta{சக} \ta{ஊழியர்} (caka ūḻiyar)}{சக ஊழியர் (caka ūḻiyar)}]{\ta{சக} \ta{ஊழியர்} (caka ūḻiyar) — a colleague}

\label{lesson:TA-C192-caka-uliyar}



\begin{quote}
Before the new one: say the Tamil for a postman, then the Tamil for a trader.
\end{quote}

\subsection*{You'll want to know: \ta{சக} \ta{ஊழியர்}}
\textbf{\ta{சக} \ta{ஊழியர்}} — \emph{caka ūḻiyar} — \textquotedblleft{}a colleague\textquotedblright{}.

\begin{grammarlens}[title={What you've built}]
One more word for this chapter.
\end{grammarlens}

\subsection*{Your turn}
\begin{itemize}
\raggedright
  \item \emph{Say it:} \emph{caka ūḻiyar}
  \item \emph{Say it:} \emph{caka ūḻiyar}, once more
  \item \emph{Say it:} say \emph{caka ūḻiyar}
  \item \emph{Recall:} say the Tamil for a postman, then the Tamil for a trader, then say \emph{caka ūḻiyar} again
\end{itemize}

\begin{tcolorbox}[breakable,colback=teal!4,colframe=teal!35!black,title={Before you move on}]
What does \ta{சக} \ta{ஊழியர்} mean? (\textquotedblleft{}A colleague\textquotedblright{}.) Say it once more.
\end{tcolorbox}
